% !TEX root = main.tex

\section{Syntax and Semantics of the \protect\NDICE Language}\label{sec:semantics}

This section provides the syntax and denotational semantics for \NDICE, a first-order functional language extended with typical probabilistic statements and a nondeterministic coin flip. 

\subsection{Syntax}

\begin{figure}
	\begin{align*}
		\tau 
		\qmorespace{\Coloneqq} &
		\BOOL 
		~\;{}|{}~\; 
		\tau \times \tau 
		%
		\qquad\qquad
		%
		v 
		\qmorespace{\Coloneqq} 
		\TRUE  
		~\;{}|{}~\; 
		\FALSE  
		~\;{}|{}~\; 
		(v, v)
		%
		\qquad\qquad
		%
		\AEXP 
		\qmorespace{\Coloneqq} 
		x  
		~\;{}|{}~\; 
		v
		\\[.25em]
		%
		e 
		\qmorespace{\Coloneqq}&
		\AEXP  
		~\;{}|{}~\; 
		(\AEXP, \AEXP) 
		~\;{}|{}~\; 
		\FST \,\AEXP  
		~\;{}|{}~\; 
		\SND \,\AEXP  
		\\
		&\quad{}~\;{}|{}~\;
		\ITE{\AEXP}{e}{e} 
		~\;{}|{}~\; 
		\LET{x = e}{e} 
		~\;{}|{}~\; 
		\FUNCNAME(\AEXP)
		\\
		&\quad{}~\;{}|{}~\;
		\FLIP{\theta} 
		~\;{}|{}~\; 
		\NFLIP 
		~\;{}|{}~\; 
		\OBSERVE{\AEXP} 
		\\[.25em]
		%
		\mathit{fun} 
		\qmorespace{\Coloneqq} &
		\FUNC{\FUNCNAME(x\colon \tau)\colon \, \tau}{e}\\[.25em]
		%
		p 
		\qmorespace{\Coloneqq} &
		\bullet e ~\;{}|{}~\; \mathit{fun}\;\, p
	\end{align*}%
	\caption{Syntax for the \NDICE language in A-normal form. Here, $\theta$ ranges over rational numbers in $[0,1]$}%
	\label{fig:syntax}%
\end{figure}%
%
%
The syntax for \NDICE is defined in \Cref{fig:syntax}.
\NDICE acts as a conservative extension of \DICE, and we thus also adopt the use of atomic expressions (\AEXP) to enforce an A-normal form \cite{flanagan1993essence}.
This will simplify the semantics and compilation rules without restricting the expressiveness of the language. 
For example, the tuple $(\FLIP{\frac{1}{3}}, \TRUE)$ can be expressed as $\LET{x = \FLIP{\frac{1}{3}}}{(x, \TRUE)}$.
Atomic expressions themselves are either values $v \in \VAL$ or variables $x \in \VARS$.
Much like \DICE, \NDICE expressions may be Booleans, tuples or tuple operations, function calls with non-recursive functions as well as the standard programming language features of sequential execution by let-statements or conditional branching through if-statements.
A \NDICE program $p$ is a series of function definitions 
followed by a single base expression $e$.
Later on, the practical implementation of \NDICE will relax the constraints imposed by the A-normal form, and will also introduce syntactic sugar for Boolean operations and bounded numbers.
We use these relaxations freely in our examples.

Additionally, \NDICE features two kinds of uncertainty:
Firstly, $\FLIP{\theta}$ models a \emph{probabilistic} coin flip which returns $\TRUE$ with probability $\theta \in [0, 1]$ and $\FALSE$ otherwise.
Secondly, $\NFLIP$ models a \emph{nondeterministic} choice, also taking values $\TRUE$ or \FALSE, but with no probability attached.
The use of these expressions is not restricted, and it is possible to interleave probabilistic and nondeterministic behavior arbitrarily.
Finally, \NDICE includes the statement \OBSERVE{e} to enable conditioning via observations. 
For purely probabilistic programs, this amounts to Bayesian reasoning, but \NDICE also features conditioning in the presence of nondeterminism.


\subsection{Semantics}

Probabilistic programs typically model a probability distribution over their output values. 
In the presence of nondeterminism, however, this is not sufficient as each resolution of nondeterminism may lead to a different probability distribution, and thus we generalize from distributions to sets of distributions.
This approach is grounded in category theory, where both probability and nondeterminism can be modeled as a monad. 
Namely, nondeterminism is modeled by the powerset monad, whereas probabilistic choice corresponds to the monad of finite distributions \cite{goy2020combining}.
Combining nondeterminism and probability thus amounts to composing these monads, which requires a distributive law between them \cite{beck1969distributive}.
%
Sadly, \citet{varacca2006distributing} showed that there is no distributive law between the monad of finite distributions and the powerset monad, making a suitable composition impossible.
This leads to difficulties for the design of programming languages:
A semantics which combines nondeterminism and probability naively will not be compositional, and the semantics of $c_1; (c_2; c_3)$ may differ from $(c_1; c_2); c_3$.

To address this, many approaches to combining the probabilistic and nondeterministic monad have been developed, for example \citet{varacca2006distributing} generalize the probabilistic monad to indexed valuations and \citet{goy2020combining} use weak distributive laws instead.
Our semantics will follow the work of \citet{mislove2000nondeterminism} and \citet{tix2009semantic}, which retains compositionality by requiring the sets of distributions to be convex.
We call a set $S$ of distributions convex, if for all distributions $d_1, d_2 \in S$ and $p \in [0,1]$, the distribution $p\cdot d_1 + (1-p)\cdot d_2$ is also in $S$.
Hence, we will use convex sets to model nondeterministic behavior, which has been applied to programming languages before \cite{zilberstein2025demonic, varacca2006distributing, jifeng1997probabilistic}. 
\\


\begin{figure}
		\begin{center}
		\begin{adjustbox}{max width=\textwidth}
				\renewcommand{\arraystretch}{2.1}
				\begin{tabular}{c@{\qquad}c}
					\multicolumn{2}{c}{ 
					% value 
					$\llbracket v \rrbracket \ttriangleq \{1|v\rangle \}$ 
					\hfill
					
					% flip
					$\llbracket \FLIP{\theta} \rrbracket \ttriangleq \{\theta|\TRUE\rangle + (1 - \theta)|\FALSE\rangle\}$
					
					\hfill
					%nflip
					$\llbracket \NFLIP \rrbracket \ttriangleq
					% semantics option 1
					\{\theta|\TRUE\rangle +  (1-\theta)|\FALSE\rangle \,|\, \theta \in [0,1]\}$

					}
					\\
					
					
					\multicolumn{2}{c}{ 
						
						% fst(v,v)
						$\llbracket \FST\,(v_1, v_2) \rrbracket \ttriangleq \{1|v_1\rangle \}$
						\hfill
						
						
						% snd(v,v)
						$\llbracket \SND\,(v_1, v_2) \rrbracket \ttriangleq \{1|v_2\rangle \}$
						\hfill
						
						% function call
						$\llbracket \FUNCNAME(v) \rrbracket^T \ttriangleq T(\FUNCNAME)(v)$
					}
					\\[1em]
					%
					% observe
					$\llbracket \OBSERVE{v} \rrbracket \ttriangleq \begin{cases}
						\{1|\TRUE\rangle\},	&\text{if } v = \TRUE\\
						\{0|\TRUE\rangle\},	&\text{if } v = \FALSE\\
						\emptyset, 		& \text{otherwise}\\
					\end{cases}$
					&
					
					%if
					$\llbracket \ITE{v_g}{e_1}{e_2} \rrbracket \ttriangleq \begin{cases}
						\llbracket e_1\rrbracket, 	& \text{if } v_g = \TRUE\\
						\llbracket e_2\rrbracket, 	& \text{if } v_g = \FALSE\\
						\emptyset,				& \text{otherwise}\\
					\end{cases} $
					\\
					
					%let
					% d_1(v) > 0 ensures that we only plug in correctly typed values, if we omit it there are problems with the observe not generating any distributions and the existential quantifiers
					\multicolumn{2}{c}{ 
						\renewcommand{\arraystretch}{1.2}
						$\begin{matrix}
							\phantom{space maker}\\
							d \iin \llbracket \LET{x = e_1}{e_2} \rrbracket\\
							\qmorespace{\Longleftrightarrow} 
							\exists\, d_1 \in \llbracket e_1 \rrbracket~{}~{}~
							\forall\, v, d_1(v) > 0~{}~{}~
							\exists\, d_v \in \llbracket e_2[x \mapsto v] \rrbracket~{}~{}~
							\forall\, v'\colon \quad
							d(v') = \sum\limits_{v,\; d_1(v) > 0 } d_1(v) \cdot d_{v}(v')
						\end{matrix}$
					}
						\end{tabular}
			\end{adjustbox}
		\end{center}
	\caption{Semantics of \NDICE expressions. The function table $T$ is left implicit in most rules.}
	\label{fig:semantics} 
\end{figure}




Before we present the semantics of \NDICE, we will briefly introduce some notations. 
Given a finite set $X$, we call a function $d \colon X \rightarrow [0,1]$ a distribution if $\sum_{x \in X} d(x) = 1$, and we call $d$ a \textit{sub}-distribution if $\sum_{x \in X} d(x) \leq 1$. 
We use bra-ket notation $d = p_1{| v_1 \rangle} + {\ldots} +  p_n {|v_n \rangle}$ to denote a sub-distribution with $d(v_i) = p_i$.
We denote the set of all sub-distributions over $X$ as $\mathcal{D}(X)$, and we denote the powerset of $X$ by $\mathcal{P}(X)$.
%
We then define the semantics of \NDICE expressions as convex sets of (sub-)distributions.
Formally, the semantics function $\llbracket \:\cdot\: \rrbracket^T\colon \NDICE \rightarrow  \mathcal{P}(\mathcal{D}(\mathit{Val}))$ associates to each \NDICE expression a set of sub-distributions over values which may be generated by it.
The semantics function is parametrized with a function table $T$ to store the behavior of functions.
Note that the function table is only needed for function calls, therefore we omit $T$ in most cases.

The semantics for closed \NDICE expressions are defined in \Cref{fig:semantics}, and they conservatively extend the semantics of \DICE~\cite{holtzen2020scaling} for well-formed programs.
For the sake of completeness, we define the semantics of ill-formed programs, e.g.\ programs which use non-Boolean values as guards of if-statements, to be the empty set.
From now on we will assume that programs are well-formed, as our inference algorithm only considers well-formed programs.

We begin by considering the semantics of values and tuple operations.
Notably, these expressions contain neither (a) nondeterminism nor (b) randomness.
Their semantics are hence (a) \emph{singleton} sets with only one (b) \emph{Dirac} distribution assigning probability 1 to a single value.
The probabilistic $\FLIP{\theta}$ is not nondeterministic but is random.
Thus, its semantics is a still a singleton but containing a non-Dirac distribution over \textit{two} possible values. 
The semantics of $\ITE{v_g}{e_1}{e_2}$ is standard. 
Due to the A-normal form, condition $v_g$ is always a value.
If it is $\TRUE$, then the if-statement behaves like $e_1$, otherwise it behaves like $e_2$. 
Next, we examine the remaining rules in greater detail.

\paragraph{Nondeterministic Flips}
The statement $\NFLIP$ is nondeterministic, but not inherently random. 
Still, we want non-determinism to act as an abstraction for all possible decision procedures we could use to return either $\TRUE$ or $\FALSE$, which includes random choices.
Thus, the semantics of $\NFLIP$ do not include just the deterministic distributions $1|\TRUE\rangle$ and $1|\FALSE\rangle$, but also all possible distributions in between.
Additionally, the semantics of $\NFLIP$ then form a convex set.

\paragraph{Let Statement}
The semantics of $\LET{x=e_1}{e_2}$ is a set of distributions that emerges from all resolutions of the nondeterminism in both $e_1$ and $e_2$.
For this, we first pick one distribution $d_1$ from the set $\llbracket e_1 \rrbracket$, thus resolving the nondeterminism in $e_1$.
Now, since there may be many values $v$ with $d_1(v) > 0$, and $e_2$ may depend on $v$, we must next pick a distribution $d_v$ from $\llbracket e_2 [x\mapsto v]\rrbracket$ -- where $e_2 [x\mapsto v]$ is $e_2$ with all free occurrences of $x$ replaced by $v$ -- for all possible intermediate values $v$.
We have now fully resolved the nondeterminism within the let statement, and the probability to output a value $v'$ can be found by considering all possible execution paths, so that $e_1$ first terminates in an intermediate value $v$, and then $e_2[x\mapsto v]$ terminates in $v'$.
Like this, we can determine the distribution for \textit{one} possible resolution to the nondeterminism. 
The semantics of the let-statement as a whole are generated by considering \textit{all} possible resolutions to the nondeterminism, which is done by considering all combinations of distributions which could be picked from the semantics of $e_1$ and $e_2$.
To illustrate this, consider the following example:
\begin{align*}
	\LET{x = \NFLIP}{\LET{y = \FLIP{\tfrac{2}{3}}}{x \leftrightarrow y}} \tag{ExLet1}\label{ExLet1}\\
%	\LET{x = \FLIP{\tfrac{2}{3}}}{\LET{y = \NFLIP}{x \leftrightarrow y}}  \tag{ExLet2}\label{ExLet2}\\
	\llbracket ExLet1 \rrbracket = \left\{p | \TRUE \rangle + (1-p)| \FALSE \rangle \,|\, p \in \left[\tfrac{1}{3},\tfrac{2}{3}\right]\right\}
\end{align*}
To generate a distribution from $\llbracket \text{\ref{ExLet1}} \rrbracket$, we first pick a distribution from $\llbracket \NFLIP \rrbracket$, for example $d_1 = \frac{1}{4}|\TRUE\rangle + \frac{3}{4}|\FALSE\rangle$.
This gives us a two values with $d_1(v) > 0$, namely $\TRUE$ and $\FALSE$.  
Thus, we have to pick distributions from the two singleton sets $\llbracket \LET{y = \FLIP{\tfrac{2}{3}}}{{\color{orange}\TRUE} \leftrightarrow y} \rrbracket = \left\{\frac{2}{3}|\TRUE\rangle + \frac{1}{3}|\FALSE\rangle\right\}$ and $\llbracket \LET{y = \FLIP{\tfrac{2}{3}}}{{ \color{orange} \FALSE} \leftrightarrow y} \rrbracket = \left\{\frac{1}{3}|\TRUE\rangle + \frac{2}{3}|\FALSE\rangle\right\}$.
Combining the three distributions then generates a possible distribution for the overall let statement, namely:
$$
\left(\tfrac{1}{4}\cdot \tfrac{2}{3} + \tfrac{3}{4}\cdot \tfrac{1}{3}\middle)\middle|\TRUE\right\rangle + 
\left(\tfrac{1}{4}\cdot \tfrac{1}{3} + \tfrac{3}{4}\cdot \tfrac{2}{3}\middle)\middle|\FALSE\right\rangle 
= \left. \tfrac{5}{12}\middle|\TRUE\right\rangle + \left. \tfrac{7}{12}\middle |\FALSE \right\rangle
\in \llbracket \text{\ref{ExLet1}} \rrbracket 
$$
%
Notably, the behavior of this expression changes if we swap the order of the flips:
%
\begin{align*}
	\LET{x = \FLIP{\tfrac{2}{3}}}{\LET{y = \NFLIP}{x \leftrightarrow y}}  \tag{ExLet2}\label{ExLet2}\\
	\llbracket ExLet2 \rrbracket = \{p | \TRUE \rangle + (1-p)| \FALSE \rangle \,|\, p \in [0,1]\}
\end{align*}
Intuitively, the additional distributions arise as $\NFLIP$ can act depending on the outcome of the previous flip.
More formally, the probabilistic flip may return two different values, which results in \textit{two} opportunities to pick a distribution from the semantics of the nondeterministic flip.
In particular, we may pick a different distribution depending on the outcome of  the probabilistic flip, and thus we could pick $1|\TRUE\rangle$ or $1|\FALSE\rangle$ to always match the previous flip or not.
This does not occur for \ref{ExLet1} as the outcome of the probabilistic flip is still unclear when the nondeterministic flip happens.
Hence, nondeterministic choices can depend on past outcomes, but they cannot predict the future.%



\paragraph{Observations and Conditioning}
The observe statement allows us to condition the behavior of expressions on the assumption that all observations succeed, i.e.\ the arguments to all observe statements were $\TRUE$.
Our semantics follow prior work and "filter out" failed observations by assigning them a probability of $0$ \cite{claret2013bayesian, fierens2015inference, borgstrom2011measure, huang2016application}.
We will illustrate this with the following example:
\begin{align*}
	\LET{x = \NFLIP}{\LET{ y = \FLIP{\tfrac{2}{3}} }{ \LET{t =\OBSERVE{x\lor y}}{y}} } \tag{ExObs}\label{ExObs}\\
	\llbracket ExObs \rrbracket = \left\{\left. \tfrac{2}{3}\middle|\TRUE \right\rangle + \left(\tfrac{1}{3} \cdot p\middle)\middle|\FALSE\right\rangle \;\middle|\; p\in \left[0, 1\right]\right\}
\end{align*}
The semantics of \ref{ExObs} can be explained as follows: With a probability of $\frac{2}{3}$ the second flip returns $\TRUE$, meaning the observation succeeds and the program returns $\TRUE$.
But, if the probabilistic flip returns $\FALSE$, then the observation only succeeds if the first flip returned $\TRUE$.
Thus, the probability to return $\FALSE$ depends on the resolution of the nondeterminism, which is represented by the parameter $p$.

Notably, $\llbracket \text{\ref{ExObs}} \rrbracket$ contains \textit{sub}-distributions like $\frac{2}{3}|\TRUE \rangle + \frac{1}{6}|\FALSE\rangle$, where the missing probability mass amounts to the likelihood for the observation to fail as both flips returned $\FALSE$.
We thus call the semantics in \Cref{fig:semantics} \textit{unnormalized semantics}. 
Retrieving the normalized semantics of an expression $e$ is straightforward, as we simply normalize the distributions in $\llbracket e \rrbracket$.
To do so, we define the normalizing constant $d_D$ of a distribution $d$ and the normalized semantics $\llbracket e \rrbracket_D$ of an expression $e$:
%
$$ d_D \coloneqq \sum_{v} d(v) \qquad \qquad \llbracket e \rrbracket_D \coloneqq \left\{d' \,\left|\; \exists d\in \llbracket e \rrbracket \colon \forall v \colon d'(v) = \frac{d(v)}{d_D} \right. \right\}$$
%
For $d_D = 0$ we define the fraction to be $0$.
The semantics $\llbracket e \rrbracket_D$ describe all output distributions of $e$, conditioned by all observations succeeding.
For example, consider $\llbracket \text{\ref{ExObs}} \rrbracket_D$ compared to $\llbracket \text{\ref{ExObs}} \rrbracket$, where $\TRUE|\frac{2}{3}\rangle + \frac{1}{3} |\FALSE \rangle$ remains unchanged, but $\frac{2}{3}|\TRUE \rangle$ will be normalized to $1|\TRUE\rangle$.
Considering unnormalized semantics first and normalizing at the end is standard, and we share this approach with \DICE and other PPLs \cite{claret2013bayesian, fierens2015inference}.


\paragraph{Function Calls and Programs}
The semantics of a function $\FUNCNAME(x \colon \tau_1)$ takes an input value $v$ of type $\tau_1$ and returns the semantics of $\FUNCNAME$'s function body $e$ (i.e.\ a set of distributions), but in which each occurrence of $x$ is replaced by $v$.
Hence, the semantics of a \NDICE function $\FUNCNAME$ is of type $\llbracket \FUNC{\FUNCNAME(x\colon \tau_1)\colon \, \tau_2}{e} \rrbracket: \mathit{Val} \rightarrow \mathcal{P}(\mathcal{D}(\mathit{Val}))$ and given formally by:
%
$$\llbracket \FUNC{\FUNCNAME(x\colon \tau_1)\colon \, \tau_2}{e} \rrbracket (v) \triangleq \llbracket e[x \mapsto v]\rrbracket$$
%
In \Cref{fig:semantics}, we use this semantical object as stored in the function table $T$ to define the semantics for function calls: 
We look up the identifier $\FUNCNAME$ of the called function in the function table $T$ and return the corresponding set of unnormalized distributions for the input value $v$.

We can now define the semantics of a \NDICE program $p$ as $\llbracket p \rrbracket^T \colon  \mathcal{P}(\mathcal{D}(\mathit{Val}))$.
We determine $\llbracket p \rrbracket^T$ by first adding each function to the function table, until eventually the base expression $\bullet \, e$ is reached.
Finally, the semantics of the program as a whole are the semantics of the base expression $e$ in the context of the accumulated function table $T$.
This behavior is formalized by the following rules, where $\bullet$ denotes the empty sequence and $\eta(fun)$ retrieves the function's identifier:
$$ \llbracket \bullet \, e \rrbracket^T \triangleq \llbracket e \rrbracket^T \qquad \llbracket fun \; p \rrbracket^T \triangleq \llbracket p \rrbracket^{T \cup \{\eta(fun) \rightarrow \llbracket fun \rrbracket^T\}} $$

